% Margin Notes Article -- a single-column article with a wide outer margin.
% Notes, citations and small figures live beside the text; the page design
% (grid, fonts, accent colour, headings) is in margin-notes.sty.
% Build with: latexmk -pdf main.tex
\documentclass[11pt]{article}
\usepackage{margin-notes}

\usepackage{amsmath}
\usepackage{pgfplots}
\pgfplotsset{compat=1.18}
\usepgfplotslibrary{groupplots,fillbetween}

% Author-year references; every \sidecite also prints the full entry in the margin.
\usepackage[backend=biber, style=authoryear, maxcitenames=2, maxbibnames=6,
  giveninits=true, useprefix=true, uniquename=false, uniquelist=false, dashed=false]{biblatex}
\addbibresource{refs.bib}
\AtEveryBibitem{\clearfield{issn}\clearfield{doi}\clearfield{url}}
\AtBeginBibliography{\renewcommand*{\mkbibnamefamily}[1]{\textsc{#1}}}
\newcommand{\sidecite}[2][]{\sidenote{\fullcite[#1]{#2}.}}

\usepackage[hidelinks]{hyperref}
\usepackage[capitalise,noabbrev]{cleveref}

% ---- Title block: edit these ------------------------------------------------
\title{Light before Dawn}
\subtitle{Street lighting advances the first song of the day in four urban songbirds}
\author{Mira Okafor, Tomás Lindqvist \& Hanne Verbeek}
\affiliation{Halden Institute for Urban Ecology, Department of Behavioural Biology}
\margininfo{\textcolor{accent}{\scshape\textls[60]{preprint}}\\
  October 2026\\[4pt]
  Correspondence to\\ \textit{m.okafor@halden-ecology.org}\\[4pt]
  \textit{Keywords}: artificial light at night, dawn chorus, acoustic monitoring, mixed models}
\shorttitle{Light before Dawn}
\shortauthors{Okafor, Lindqvist \& Verbeek}

\begin{document}
\maketitle

\epigraph{Before the light came, the blackbird had begun; it sang as though the
  lamp outside the window were the first edge of morning.}
  {field notebook, 14 April 2024}

\begin{abstract}
Songbirds time their first song of the day against the rising light, and the
city adds light of its own. We analysed 4\,545 species-mornings at 38 sites along a
measured gradient of street lighting and estimated, for four common species, how
far the onset of song moves for each tenfold increase in ground-level
illuminance. Robins and blackbirds sing more than ten minutes earlier per decade
of light, great tits about four, and chaffinches barely at all. The ordering
follows eye size, which suggests that the response is set by how much light a
species needs to see, not by its place in the chorus.
\end{abstract}

\section{Introduction}

The dawn chorus is among the most regular events in temperate ecology. Each
species begins at a characteristic time before sunrise, and the order of
entries is stable enough that field ornithologists once used it as a clock.%
\sidecite{marler2004} The cue behind that regularity is light: as the sky
brightens through civil twilight, each species starts once illuminance crosses
a threshold of its own.

Cities disturb the cue. A sodium or LED street lamp delivers several lux to a
nearby perch through the whole night, which is brighter than the sky at the
start of civil twilight. The best-known study of the effect found that males
of several species close to street lights began singing earlier and fathered
more extra-pair young \parencite{kempenaers2010}. Later work has shown that the
spectrum of the lamp matters as much as its intensity\sidecite{hale2017} and
that the size of a species' eye predicts how strongly it reacts.%
\sidecite{moreau2022}

Most of this evidence comes from a handful of sites, compared as lit against
unlit. Light at night is not a binary treatment, however; it falls off with
distance and is shaded by trees and walls.\footnote{A lamp ten metres away
  delivers roughly a hundredth of the illuminance of one directly above the
  perch, so two territories on the same street can differ by two orders of
  magnitude.} We therefore treat light as a continuous dose and ask a simpler
question: how many minutes earlier does each species start for each tenfold
increase in light?

\section{Methods}

\subsection{Sites and light measurements}

We placed autonomous recorders at 38 sites across a city of about 300\,000
people and its rural edge, chosen to span five orders of magnitude of
night-time illuminance. At each site we measured horizontal illuminance at
perch height with a calibrated lux meter on three moonless nights, following an
established ground protocol,\sidecite{dacosta2019} and used the median as the
site value~$L_i$. Recording ran from 1 March to 31 May in two seasons.

\subsection{Defining the onset of song}

\begin{marginfigure}
\begin{tikzpicture}
\begin{axis}[
    width=56mm, height=40mm, scale only axis=false,
    axis lines=left, xmin=-90, xmax=10, ymin=0, ymax=1.08,
    xtick={-80,-40,0}, ytick={0,0.25,1},
    yticklabels={0,$\kappa$,1},
    tick label style={font=\scriptsize}, label style={font=\scriptsize},
    xlabel={minutes from civil dawn}, ylabel={$r(t)/r_{\max}$},
    ylabel style={yshift=-6pt}, xlabel style={yshift=3pt},
    tick align=outside, tick style={black!60}, axis line style={black!60},
  ]
  \addplot[domain=-90:10, samples=80, thick, accent] {1/(1+exp(-(x+38)/7))};
  \draw[dashed, black!45] (axis cs:-90,0.25) -- (axis cs:10,0.25);
  \draw[black!60] (axis cs:-45.7,0) -- (axis cs:-45.7,0.25);
  \fill[accent] (axis cs:-45.7,0.25) circle (1.3pt);
  \node[font=\scriptsize, anchor=south east] at (axis cs:-46,0.28) {$t^{\ast}$};
\end{axis}
\end{tikzpicture}
\caption{Onset $t^{\ast}$ is the first minute at which the smoothed song rate
  reaches a fraction $\kappa$ of its morning maximum.}
\label{fig:onset}
\end{marginfigure}
For each species, morning and site we counted songs per minute with a
classifier trained on our own annotated recordings\sidecite{verbeek2021} and
smoothed the counts with a five-minute running mean $r(t)$. The onset of song
is the first time at which this rate reaches a fixed fraction of its maximum
for that morning (\cref{fig:onset}):
\begin{equation}
  t^{\ast} = \min\bigl\{\, t : r(t) \ge \kappa \, r_{\max} \bigr\},
  \qquad \kappa = 0.25 .
  \label{eq:onset}
\end{equation}
A threshold relative to the morning maximum makes the onset insensitive to
how far a bird sings from the recorder. We express $t^{\ast}$ in minutes
relative to the start of civil twilight, so that negative values mean song
before civil dawn.

\subsection{Model}

We fitted one linear mixed model per species $j$, with onset for morning $i$
explained by the logarithm of site illuminance, day of the season $D_i$ and a
random intercept for site $s[i]$:
\begin{equation}
  t^{\ast}_{ij} = \alpha_j + \beta_j \log_{10}\!\left(L_{s[i]} + 1\right)
    + \gamma_j D_i + u_{s[i]} + \varepsilon_{ij},
  \qquad u_s \sim \mathcal{N}\!\left(0, \sigma_{u}^{2}\right).
  \label{eq:model}
\end{equation}
The slope $\beta_j$ is the quantity of interest: the shift in onset, in
minutes, for each tenfold increase in light. Adding one lux before the
logarithm keeps unlit sites finite. Models were fitted in
R,\sidecite{rcore2025} and intervals are 95\% profile likelihood intervals.

\section{Results}

\begin{table}[tb]
\begin{sidecaptioned}{\caption{Shift in onset of song per tenfold increase in
    illuminance, from \cref{eq:model}. Negative values are earlier song. The
    last column is the predicted shift at 30~lux, a bright residential
    street.}\label{tab:slopes}}
\centering\small
\begin{tabular}{@{}l r r r@{\hspace{0.6em}}l r@{}}
  \toprule
  Species & Sites & Mornings & \multicolumn{2}{c}{$\beta_j$ (min)} & 30 lux \\
  \midrule
  Robin      & 34 & 1\,288 & $-11.4$ & [$-13.1$, $-9.7$] & $-17.0$ \\
  Blackbird    & 36 & 1\,344 & $-9.2$  & [$-10.6$, $-7.8$] & $-13.7$ \\
  Great tit     & 31 & 1\,107 & $-4.1$  & [$-5.3$, $-2.9$]  & $-6.1$  \\
  Chaffinch   & 22 &    806 & $-1.3$  & [$-2.6$, $0.0$]   & $-1.9$  \\
  \bottomrule
\end{tabular}
\end{sidecaptioned}
\end{table}

All four species began singing earlier at brighter sites, but by very
different amounts (\cref{tab:slopes}). Robins responded most strongly: at a
site lit to 30~lux, a typical bright residential street, they started about
seventeen minutes before robins at dark sites. Blackbirds were close behind.
Great tits shifted by less than half as much, and the chaffinch slope was
indistinguishable from zero.

The relation is close to linear on the logarithmic scale across the whole
gradient (\cref{fig:gradient}). We found no sign of a threshold below which
light has no effect: even the dimmest lit sites, at a fraction of a lux, sat
above the dark-site mean for robins and blackbirds.\footnote{A model with a
  breakpoint in $\log_{10} L$ did not improve on the linear one for any
  species ($\Delta\mathrm{AIC} < 2$).}

\begin{widefigure}[t]
\centering
\begin{tikzpicture}
\pgfmathsetseed{2026}
\begin{groupplot}[
    group style={group size=4 by 1, horizontal sep=7mm, yticklabels at=edge left},
    width=46mm, height=44mm, scale only axis=false,
    xmin=-0.05, xmax=1.65, ymin=-82, ymax=8,
    xtick={0,0.5,1,1.5}, xticklabels={0,2,9,31},
    ytick={-80,-60,-40,-20,0},
    axis lines=left, tick align=outside,
    tick style={black!55}, axis line style={black!55},
    tick label style={font=\scriptsize}, label style={font=\scriptsize},
    title style={font=\footnotesize\itshape, yshift=-4pt},
    xlabel={illuminance (lux)},
  ]
  \nextgroupplot[title={Robin}, ylabel={onset (min from civil dawn)}]
    \addplot[only marks, mark=*, mark size=0.9pt, black!40, samples=34, domain=0:1.6] {-52 - 11.4*x + 7*rand};
    \addplot[name path=a1, draw=none, domain=0:1.6] {-52 - 13.1*x + 2.5};
    \addplot[name path=b1, draw=none, domain=0:1.6] {-52 - 9.7*x - 2.5};
    \addplot[accent!18] fill between[of=a1 and b1];
    \addplot[accent, thick, domain=0:1.6] {-52 - 11.4*x};
  \nextgroupplot[title={Blackbird}]
    \addplot[only marks, mark=*, mark size=0.9pt, black!40, samples=36, domain=0:1.6] {-44 - 9.2*x + 7*rand};
    \addplot[name path=a2, draw=none, domain=0:1.6] {-44 - 10.6*x + 2.3};
    \addplot[name path=b2, draw=none, domain=0:1.6] {-44 - 7.8*x - 2.3};
    \addplot[accent!18] fill between[of=a2 and b2];
    \addplot[accent, thick, domain=0:1.6] {-44 - 9.2*x};
  \nextgroupplot[title={Great tit}]
    \addplot[only marks, mark=*, mark size=0.9pt, black!40, samples=31, domain=0:1.6] {-31 - 4.1*x + 6*rand};
    \addplot[name path=a3, draw=none, domain=0:1.6] {-31 - 5.3*x + 2.2};
    \addplot[name path=b3, draw=none, domain=0:1.6] {-31 - 2.9*x - 2.2};
    \addplot[accent!18] fill between[of=a3 and b3];
    \addplot[accent, thick, domain=0:1.6] {-31 - 4.1*x};
  \nextgroupplot[title={Chaffinch}]
    \addplot[only marks, mark=*, mark size=0.9pt, black!40, samples=22, domain=0:1.6] {-12 - 1.3*x + 6*rand};
    \addplot[name path=a4, draw=none, domain=0:1.6] {-12 - 2.6*x + 2.6};
    \addplot[name path=b4, draw=none, domain=0:1.6] {-12 - 0.0*x - 2.6};
    \addplot[accent!18] fill between[of=a4 and b4];
    \addplot[accent, thick, domain=0:1.6] {-12 - 1.3*x};
\end{groupplot}
\end{tikzpicture}
\caption{Site-mean onset of song against night-time illuminance, on the
  logarithmic scale of \cref{eq:model}. Lines are fitted slopes with 95\%
  intervals; grey points are sites. The same axis is used for all species, so
  the panels read from left to right in the order of the chorus.}
\label{fig:gradient}
\end{widefigure}

\paragraph{Season} Onset advanced through the spring at a similar rate in all
species, by 0.4 to 0.6 minutes a day, as civil dawn itself moved earlier. The
light effect did not change with day of the season for any species, so the
shift is not confined to the start of breeding.

\paragraph{Spectrum} Fourteen sites were converted from sodium to white LED
lamps between the two seasons. Holding illuminance fixed, robins at converted
sites started a further 3.8 minutes earlier [1.2, 6.4], consistent with the
greater sensitivity of passerine vision to short wavelengths.%
\sidecite{hale2017}

\section{Discussion}

The ranking of the four species matches the ranking of their eye size
relative to body mass, which is largest in robins and smallest in chaffinches.%
\sidecite{moreau2022} A large eye collects more light, so a species with a
large eye sees well enough to forage and defend a territory earlier in
twilight, and a street lamp moves that moment earlier still. The response
would then be a matter of vision rather than of position in the chorus, which
is the explanation favoured in a broad review of night-light impacts
\parencite{gaston2013}.

Two cautions apply. Light and noise are correlated in cities, and although we
measured traffic noise at every site, the two cannot be fully separated in an
observational design. And an earlier start need not mean a longer day of song:
whether lit birds also stop earlier in the evening is a question for the
evening recordings, which we have not yet analysed.

The practical message is simple. Because the effect is logarithmic, the first
lux matters most: dimming a bright street from 30 to 10~lux recovers less of
the dark-site timing than shielding a dark one from its first stray lamp.

\paragraph{Acknowledgements} We thank the residents who hosted recorders in
their gardens, and the city lighting department for switching schedules.

\printbibliography[heading=bibintoc, title={References}]

\end{document}
